\documentclass{CSUniSchoolLabReport}

\usepackage[italian]{babel}
\usepackage{fancyhdr}
\graphicspath{ {./images/} }


\title{GreenTrails - Modulo AI}
\author{\textsc{Roberta Galluzzo}, \textsc{Gabriele Di Stefano}}

\pagestyle{fancy}
\renewcommand{\headrulewidth}{0pt}
\fancyhf{}
\lfoot{\includegraphics[scale=0.05]{logo}
		C03 - Modulo AI}
\rfoot{\thepage}

\begin{document}
\maketitle

\begin{center}
	\includegraphics[scale=0.2]{logo}
\end{center}

\begin{center}
	\begin{tabular}{l r}
		Team: & C03 GreenTrails \\
		Progetto: & Progetto combinato IS/FIA 2023-2024
	\end{tabular}
\end{center}

\pagebreak

\tableofcontents

\pagebreak

\section{Definizione del problema}
\subsection{Introduzione}

Negli ultimi anni, l'ecosostenibilità è diventata una priorità, poiché il mondo, in rapida crescita, deve prestare maggiore attenzione all'ambiente. In questo contesto, l'applicazione web "GreenTrails" si rivolge al settore turistico, con l'obiettivo di facilitare l'organizzazione di viaggi sostenibili, tenendo conto delle sfide che gli utenti affrontano in questo ambito.
L'obiettivo di "GreenTrails" è permettere agli utenti di pianificare facilmente e rapidamente viaggi ecosostenibili, aiutandoli a individuare mete, strutture e attività adatte. \\\\
L'applicazione, dunque, mira a:
\begin{itemize}
    \item Facilitare la ricerca di strutture eco-friendly;
    \item Generare automaticamente itinerari in base alle preferenze personali degli utenti;
    \item Velocizzare il contatto con le strutture selezionate.
\end{itemize}


\subsection{Obiettivi}

GreenTrails prevede, appunto, la \textbf{generazione automatica di itinerari multi-attività in base alle preferenze specificate dall'utente}, attraverso l'utilizzo di un algoritmo di intelligenza artificiale. Le preferenze sono così composte:
\begin{itemize}
    \item \textbf{Località preferita}: una scelta tra \textit{Mare}, \textit{Montagna}, \textit{Città}, oppure \textit{Nessuna preferenza}.
    \item \textbf{Tipologia di alloggio preferita}: uno tra \textit{Hotel}, \textit{Bed \& Breakfast}, \textit{Villaggio turistico}, \textit{Ostello} oppure \textit{Nessuna preferenza}.
    \item \textbf{Preferenze alimentari}: \textit{Vegan}, \textit{Vegetarian}, \textit{Gluten-free} o \textit{Nessuna preferenza}.
    \item \textbf{Tipologia di attività preferita}: una categoria a scelta tra \textit{All'aperto}, \textit{Visite storico-culturali}, \textit{Gastronomia} oppure \textit{Nessuna preferenza}.
    \item \textbf{Animale domestico}: esprime la preferenza di portare con sé un animale da compagnia, con una semplice scelta tra \textit{Sì} oppure \textit{No}.
    \item \textbf{Budget}: \textit{Basso}, \textit{Medio}, \textit{Alto} oppure \textit{Flessibile}.
    \item \textbf{Souvenir locali}: esprime l'interesse nell'acquisto di prodotti artigianali locali, con una semplice scelta tra \textit{Sì} oppure \textit{No}.
    \item \textbf{Stagione preferita}: \textit{Autunno-inverno}, \textit{Primavera-estate} o \textit{Nessuna preferenza}.
\end{itemize}
~\\
Inoltre, tale algoritmo dovrà rispettare il vincolo di dover generare un percorso valido entro il tempo massimo prestabilito dal sistema (\textbf{2 secondi}, come definito dal System Design Goal \textit{"DG\_12 - Tempo di risposta"}).


\subsubsection{Ulteriori considerazioni}

In mancanza di ulteriori specifiche, che avrebbero potuto maggiormente impostare il problema, si è deciso di imporre dei vincoli aggiuntivi per facilitare la fase di formulazione:
\begin{itemize}
	\item il modulo di intelligenza artificiale non cercherà su tutto il database dell'applicazione, ma soltanto su un subset fornito dal sistema stesso (in altre parole, \textbf{consideriamo soltanto le attività di una città o una specifica area} fornite dal sistema);
	\item il percorso dovrà \textbf{terminare con una e una sola struttura ricettiva};
	\item il percorso dovrà \textbf{contenere almeno due attività turistiche, fino ad un massimo di quattro};
	\item le attività del percorso possono essere \textbf{distanti tra loro al più 5 km}.
\end{itemize}

\pagebreak

\subsection{Specifica PEAS}

\begin{itemize}
 \item \textbf{Performance}: la misura di prestazione dell'agente, in questo caso, si basa sulla pertinenza del percorso fornito con le preferenze dell'utente e sulla sua distanza complessiva.
 \item \textbf{Environment}: l'ambiente di riferimento è l'applicazione web di GreenTrails, la quale comunica direttamente con l'agente.
 \item \textbf{Actuators}: sono le soluzioni restituite all'applicazione web.
 \item \textbf{Sensors}: sono i dati in input (ovvero, la lista delle attività disponibili e le preferenze dell'utente) forniti dall'applicazione web.
\end{itemize}


\subsection{Caratteristiche dell'ambiente}

L'ambiente in cui opera l'agente ha le seguenti caratteristiche:
\begin{itemize}
 \item è \textbf{completamente osservabile}, in quanto prende in considerazione tutte le possibili attività fornite dal sistema;
 \item è \textbf{stocastico}, poiché lo stato varia indipendentemente dalle azioni dell'agente;
 \item è \textbf{episodico}, dal momento che la soluzione scelta dall'agente varia in base alla lista di attività e alle preferenze dell'utente (dunque, ad ogni episodio);
 \item è \textbf{statico}, in quanto non subisce variazioni mentre l'agente sta valutando le possibili soluzioni;
 \item è \textbf{discreto}, perché è possibile identificare un numero finito di scelte (tutte le possibili combinazioni di attività e strutture ricettive);
 \item è \textbf{ad agente singolo}, dato che consente la presenza di un unico agente.
\end{itemize}

\pagebreak

\subsection{Analisi del problema}

Considerando la natura del problema, che è di \textbf{ottimizzazione}, possiamo optare per l'utilizzo di un \textbf{algoritmo genetico}. Tale scelta è condizionata anche dalla \textbf{necessità progettuale} (per il progetto di Ingegneria del Software) di utilizzare un algoritmo genetico per la generazione degli itinerari. \\
Formalizzando:
\begin{itemize}
    \item \textbf{Spazio degli stati}: tutte le possibili permutazioni delle attività contenenti almeno due attività turistiche ed esattamente una struttura ricettiva al termine, distanti tra loro al più 5 km;
    \item \textbf{Stato iniziale}: percorso vuoto;
    \item \textbf{Azioni}: aggiungere un'attività turistica (fino ad un massimo di 4), o una struttura ricettiva (se è al termine del percorso);
    \item \textbf{Modello di transizione}: restituisce un percorso contenente l'attività aggiunta alla fine dello stesso;
    \item \textbf{Test obiettivo}: il percorso deve contenere attività basate sulle preferenze utente, e deve essere realizzato entro 2 secondi;
    \item \textbf{Costo di cammino}: distanza dall'attività successiva, in km.
\end{itemize}

\section{Soluzione del problema}
\subsection{Tecnologie utilizzate}

Si è scelto di utilizzare Python come linguaggio di programmazione. Le librerie utilizzate, invece, includono FastAPI (per la realizzazione di un'API HTTP semplice).

\subsection{Set di dati utilizzati}

Per ciò che concerne il reperimento dei dati, sono disponibili due opzioni:

\begin{itemize}
    \item la prima, che consiste nell'\textbf{utilizzare un dataset pubblicamente disponibile};
    \item la seconda, che fa uso di un \textbf{Large Language Model per la generazione di dati artificiali}.
\end{itemize}

Abbiamo optato per la generazione di dati artificiali, in quanto è irrilevante avere dati veritieri in questo contesto: \textbf{il focus del problema è quello di trovare dei percorsi affini in base alle esigenze dell'utente}, ed è perciò non necessario utilizzare informazioni di strutture e utenti reali. \\
Ci siamo avvalsi di \textit{GPT-3.5} per la generazione dei dataset di visitatori e attività.

\subsubsection{Formattazione dei dati}

\subsection{Specifica dell'algoritmo genetico}

\subsubsection{Codifica degli individui}

La nostra decisione è stata quella di determinare ogni individuo come un array a taglia fissa, di dimensione 5, contenente gli ID numerici delle attività, in cui l'ultimo elemento è sempre una struttura ricettiva. Per rappresentare eventuali percorsi con un numero di attività turistiche inferiori a 4, si è inoltre deciso di utilizzare un valore sentinella al posto dell'ID, con valore 0.

\subsubsection{Inizializzazione}

Per generare una popolazione iniziale con soluzioni accettabili, si è pensato di utilizzare il seguente algoritmo:

\texttt{1. Separa l'array delle attività in un array di sole strutture ricettive e un array di sole attività turistiche. \\
2. Costruisci un K-D Tree contenente le coordinate delle sole attività turistiche. \\
3. Prendi casualmente una struttura ricettiva e rimuovila dall'array. \\
4. Estrai dal K-D Tree le 4 attività turistiche più vicine alla struttura ricettiva. \\
5a. Se le prime 2 più vicine non rispettano il vincolo di distanza di 5 km, scarta la struttura ricettiva e torna al punto 3. \\
5b. Altrimenti, considera le 2 attività turistiche rimanenti: se rispettano il vincolo, allora faranno parte dell'individuo, altrimenti verranno prese in considerazione solo quelle che lo rispettano. \\
6. Crea un individuo con la struttura ricettiva alla fine, aggiungendo a ritroso le attività in ordine di distanza; quelle non prese in considerazione avranno valore 0. \\
7. Ritorna al passo 3 finché il numero di individui è minore di n. 
}

Nell'algoritmo sopra specificato, si consideri \textit{n} la dimensione della popolazione.

\subsubsection{Selezione}

L'algoritmo di selezione considerato è il Roulette Wheel con Elitismo. Esso consiste nel preservare k individui promettenti e avviare una selezione casuale con probabilità per i restanti n-k individui. In questo modo, garantiamo che ogni generazione non sia peggiore della precedente, esponendoci, però, maggiormente al rischio di convergenza prematura.

\subsubsection{Crossover}

La strategia adottata è il Single Point Crossover, dividendo l'array al terzo elemento: in questo modo, avremo che i primi 2 elementi proverranno dal primo genitore, mentre i restanti 3 dal secondo. 

\subsubsection{Mutazione}

L'approccio utilizzato per la mutazione è lo Swap, che consiste nello scambio di due attività (ad eccezione della struttura ricettiva), scegliendo in modo casuale gli indici dell'array, a patto che la soluzione rimanga ammissibile.

\subsubsection{Funzione di fitness}



\subsubsection{Stopping condition}

\end{document}
